\documentclass{article}
\usepackage[utf8]{inputenc}
\usepackage{graphicx}
\usepackage{hyperref}

\title{It-sikkerhed: A4}
\author{Silja Maria(dfv380) og Christian Jacobsen(wbr220) }
\date{October 2020}

\begin{document}

\maketitle

\section{Resubmission}
For this resubmission, we have completed task 5 in the seed labs, and overall fixed
the rest of the report up. We also answer questions 20.7 and 20.17 correctly this
time. \\


\section{Review Questions}

\subsection{20.4}
In this question, we wish to explain what the between a block cipher and a stream cipher is.\\ 

Block cipher and stream cipher are both methods for symmetric encryption. They differ primarily in the size of the data, which they apply the encryption to. Using a block cipher method, the data will first be grouped into blocks and then each of these blocks will be encrypted by an encrypting algorithm. Using a stream cipher method on the other hand, the encryption algorithm will be applied to every byte in the data. Thus, cipher streams performs much more encryptions operations, but on a smaller input data size than cipher blocks does. 


\subsection{20.11}
In this question, we wish to explain what the difference between a session key and a master key is.\\

Both a session key and a master key are types of symmetric encryption keys, but have different purposes and different duration. A session key is used encrypted messages, when two end hosts wants to communicate. The key will then be used encrypt all messages send between the two hosts and when the session is stops, the key will be destroyed. A master key is used between a key distributor and a host and have two purposes. It can either be used to encrypt other keys or it can be used to generate other keys. An example of this is when we wish to generate a session key for two hosts. In this case, the master key will be used to protect the session key, when it is send to the two hosts. 

\subsection{21.5}
In this question, we wish to briefly explain Diffie-Hellman key exchange.\\

Diffie-Hellmans key exhange is a way to generate a session key between two users that want use encryption when communicate with each other. Before generating the session key, the two users need to agree on two numbers p and q, where p is a prime number and q is primitive root of p. Then both of the two hosts generate a secret randomly value, a and b. These values is then used to calculate a shared value of each of the hosts. These two shared values are called A and B and are calculated in this way\\

\[A = q^a \texttt{mod} p\]
\[B = q^b \texttt{mod} p\]

The two hosts now exhange these values and use them to calculate the secret key, which will be used to encrypt the messages they send to each other. The secret is calculated of both of the hosts and result in the same number:
\[K = B^a \texttt{mod} p\]
\[K = A^b \texttt{mod} p\]

Since the hosts are the only ones who knows the secret values, they will also be they only ones to know the value of the secret key. 

\subsection{22.10} 
In this question, we wish to explain the transport and tunnel modes of ESP.\\

When using ESP, the protocol can either operate in transport mode or in tunnel mode. If the protocol operates in transport mode, the protocol will only encrypt the data part in the packet. Using tunnel mode instead, the entire packet is encrypted. 

\subsection{23.6}
In this question, we wish to answer what the role of a CA in X.509 is?\\

CA is a third party that is used to control the validation of a public key. By this we mean to check if a public key belongs to the user, which it claims to belong to. Users can obtain a certificate from this third for their public key. This certificate can then verify the public key, which other users can use as a indicator to trust this key. 

\subsection{23.10}
In this question, we wish to answer how do most current X.509 implementations check the validity of signatures on a certificate?\\

Browsers, and the like, include a list of trusted Certificate Authorities, which are
responsible for signing these certificates. For example there exists the widely
used Lets Encrypt service, designed around making ssl and https connections available
for everyone. If a website has their certificate signed by Lets Encrypt, then it
can be trusted as much as the CAs signing process can be trusted.\\
It is also necessary to maintain a Certificate Revocation List(CRL) which contains
information about which certificates are revoked, and therefore shouldn't be
considered valid any longer. The book however notes, that these lists are often not
checked because of the overhead required
%hvordan tjekker x.509 om signaturen på certifcatet er rigtigt 

%x.509 har en certificate revocation list. hver certificate entry indeholder et serial nummer for en certificat 

\section{Problems}
\subsection{20.7}
In this problem, we wish to answer if we can perform encryption operations in parallel on multiple blocks of plaintext in any of the five modes? How about decryption?
\\
\\
\textbf{ECB} - using this approach, the data will be grouped into n blocks and then be encrypted using the same key. The blocks will then be put together into the total cipher text. When decrypting the cipher text, the same key will be used to decrypt all of the n blocks. Since there exists no dependency between each of the blocks when encrypted or decrypted, these operations can be performed in parallel.  
\\
\\
\textbf{CBC, CFB and OFB} - using each of these approaches, the encryption of one block will depend on the cipher text of the previous block. Thus, this cannot be performed in parallel. The applies to decryption using OFB, since the decryption of one block append on the previous block. However, using CBC or CFB the decryption can be done in parallel, since the calculation only uses the cipher text and thus, none of the calculation of the previous blocks. 
\\
\\
\textbf{CTR} - using this approach, we choose a counter value, which will be used to encrypt the data. The counter value will be increased by one for each of the blocks. The encryption of each block can be performed in parallel as long as the know the number of the block. The same with decryption. 



\subsection{20.17}
In this problem, we wish to answer the following question:\\

\textit{Suppose someone suggests the following way to confirm that the two of you are both in possession of the same secret key. You create a random bit string the length of the key, XOR it with the key, and send the result over the channel. Your partner XORs the incoming block with the key (which should be the same as your key) and sends it back. You check, and if what you receive is your original random string, you have  verified that your partner has the same secret key, yet neither of you has ever transmitted the key. Is there a flaw in this scheme?}
\\
\\ 
We will be able to use the method to verify that both parts have the same key, since XOR'ing a bit string with another bit string twice will result in the original bit string. However, the method is not very safe. If an attacker captures first the encrypted string, and then captures the original string, he/she is able to calculate the key by XOR'ing the two. 

\subsection{21.7}
In a public-key system using RSA, you intercept the ciphertext C = 61 sent to a user
whose public key is e = 11, n = 91. What is the plaintext M?\\
We use prime factorization on $n$: $n = 7\cdot 13$. We calculate
$\phi (n) = (p-1)(q-1) = 6 \cdot 12 = 72$. We know $e = 11$. We determine
$de \% 72 = 1,\quad d < 72,\quad d = 59$\\
Knowing the private key, we present the decryption of the given ciphertext:
\[
    M = C^d \% n = 61^59 \% 7 = 3
\]


%22.2
\subsection{22.2}
Consider the following threats to Web security and describe how each is countered by
a particular feature of SSL:
\subsubsection{Man-in-the-middle attack: An attacker interposes during key exchange,
acting as the client to the server and as the server to the client.}
This kind of attack is thwarted by the certificate and chain of trust ingrained in
the SSL protocols. If an attacker sets up between the server and the client, the
client will quickly realise that the certificate used for the domain isn't signed
by anyone in its trust store.
\subsubsection{Password sniffing: Passwords in HTTP or other application traffic are 
eavesdropped.}
Password-sniffing is prevented by encrypting requests before sending them out. Thus
noone without the key-pairs generated in the handshakes, will not be able to read
the sniffed packets.
\subsubsection{IP spoofing: Uses forged IP addresses to fool a host into accepting bogus 
data.}
While a connection might be spoofed, the SSL protocol demands encrypted communication,
and if a session is hijacked, the host will notice, and either abort the session or
try to renegotiate a session.
\subsubsection{IP hijacking: An active, authenticated connection between two hosts is 
disrupted and the attacker takes the place of one of the hosts.}
The attacker will in this case not know the MAC shared in the handshake, and even
if they recieve/send packets to the server/client, the recieving party will instantly
know that the session has been broken.
\subsubsection{SYN flooding: An attacker sends TCP SYN messages to request a 
connection 
but does not respond to the final message to establish the connection fully. The 
attacked TCP module typically leaves the “half-open connection” around for a few 
minutes. Repeated SYN messages can clog the TCP module.}
SSL has no direct feature to prevent this kind of flooding, especially since it sits
on top of the tcp layer.


\section{Seed labs}
For this lab session, we're working with encryption.\\
\subsection{Task 2}
We use the "openssl-enc" tool which encrypts/decrypts text using symmetric
key pairs. We try three different ciphers:
\begin{verbatim}
seed@VM:~$ cat plain.txt
Hello world

seed@VM:~$ openssl enc -aes-256-cbc -in plain.txt -out cipher.bin \ 
    -K 0011223344556677889aabbccddeeff \ 
    -iv 0 102030405060708
seed@VM:~$ xxd cipher.bin
00000000: 9b69 d212 5c5f cac4 6942 7214 dab8 5522  .i..\_..iBr...U"
seed@VM:~$ openssl enc -des3 -in plain.txt -out cipher.bin \ 
    -K 0011223344556677889aabbccddeeff -iv 0102030405060708
seed@VM:~$ xxd cipher.bin
00000000: 61c4 c596 3e2b dce8 c02b b7f5 e26c f470  a...>+...+...l.p
seed@VM:~$ openssl enc -aes-256-cbc -in plain.txt -out cipher.bin \ 
    -pass pass:encryption_is_fun
seed@VM:~$ xxd cipher.bin
00000000: 5361 6c74 6564 5f5f 4d88 db81 3e6f d9af  Salted__M...>o..
00000010: 001d eb05 b943 4767 8892 7abc 96e2 7979  .....CGg..z...yy
\end{verbatim}
First we encode using the aes-256-cbc cipher, then the des3 cipher, and finally
i show using a passphrase from which the keys and ivs can be generated, allowing for
easier decryption at a later point.

\subsection{Task 3}
Task 3 of this lab, has us encrypting a simple photo using the ECB (Electronic Code
Book) and CBC (Cipher Block Chaining) modes. As seen by the picture in 
figure \ref{fig:pic_original}, the chain-block mode is much better at obscuring the
image, while the ecb mode still leaves a noticable outline of the shapes. This is
what we expected with the chain-block depending on each previous block for randomness
as opposed to the ecb only having the current block to use.\\

It should be noted that this wouldn't be the same for any format. In a bitmap the
pixel values are encoded one after the other, and the colors in the given picture
are very uniform. When the ECB looks through the image, each equivalent block of the
image will be encrypted the same, and the bitmap will thus look like the original. If
one took a photo and encrypted it as a jpg, the image would be hidden much better,
since there aren't as many repeating blocks of pixels. To test this, we demonstrate
with a picture with both large sets of randomness, and more uniform data - A subset
of the (colored)mandelbrot set.\footnote{Courtesy of
\url{https://en.wikipedia.org/wiki/Mandelbrot_set}} The reason this has been chosen
is because it represents both large "random" sets of data, like a photograph would,
and large uniform sets of data, like our previous example. The original image is
seen on figure~\ref{fig:mandelbrot}, and a side-by-side can be seen on figure~\ref{fig:mandelbrot_enc}.\\
As should be obvious the large contiguous areas of color exhibit the previously 
defined behaviour, \emph{but} the more complex structures, arms and sattelites, show
the randomness that is to be expected of e.g. a photo.
The cbc of course is still completely unable to be read, and should be considered
boring.

\begin{figure}[!htb]
    \centering
    \includegraphics[width=\linewidth]{comparison.png}
    \caption{The side-by-side comparison. Left is the chain-block mode, and right is 
    the electronic code block}
    \label{fig:pic_original}
\end{figure}
\begin{figure}[!h]
    \centering
    \includegraphics[width=0.8\linewidth]{mandelbrot.png}
    \caption{An image showing both large random blobs and large uniform color}
    \label{fig:mandelbrot}
\end{figure}
\begin{figure}[!h]
    \centering
    \includegraphics[width=0.45\linewidth]{mandelbrot_cbc.png}%
    \includegraphics[width=0.45\linewidth]{mandelbrot_ecb.png}%
    \caption{Another side-by-side, this time of mandelbrot, with CBC on the left and 
    ECB on the right.}
    \label{fig:mandelbrot_enc}
\end{figure}
This concludes task 3\\
%\subsection{Task 4}
%For task 4 we encrypt and decrypt files to analyze their padding. This is done
%by creating files of size 5, 10 and 16, encrypting them, then decrypting them with
%the -nopad option. See below:
%\begin{verbatim}
%seed@VM:~$ echo -n "12345" > f1.txt
%seed@VM:~$ echo -n "1234567890" > f2.txt
%seed@VM:~$ echo -n "123456789abcdef0" > f3.txt
%seed@VM:~$ wc f3.txt
% 0  1 16 f3.txt
%seed@VM:~$ openssl enc -aes-128-cbc -e -in f1.txt -out f1.crypt
%seed@VM:~$ openssl enc -aes-128-cbc -e -in f2.txt -out f2.crypt
%seed@VM:~$ openssl enc -aes-128-cbc -e -in f3.txt -out f3.crypt
%seed@VM:~$ openssl enc -aes-128-cbc -d -nopad -in f1.crypt -out f1.decrypt
%seed@VM:~$ openssl enc -aes-128-cbc -d -nopad -in f2.crypt -out f2.decrypt
%seed@VM:~$ openssl enc -aes-128-cbc -d -nopad -in f3.crypt -out f3.decrypt
%seed@VM:~$ xxd f1.decrypt
%00000000: 3132 3334 350b 0b0b 0b0b 0b0b 0b0b 0b0b  12345...........
%seed@VM:~$ xxd f2.decrypt
%00000000: 3132 3334 3536 3738 3930 0606 0606 0606  1234567890......
%seed@VM:~$ xxd f3.decrypt
%00000000: 3132 3334 3536 3738 3961 6263 6465 6630  123456789abcdef0
%00000010: 1010 1010 1010 1010 1010 1010 1010 1010  ................
%\end{verbatim}
%Observed is the file sizes: 32bytes for the f1 and f2, and 48 for the f3 file.
%This is because the encrypted files get padded. Decrypting these files, we see
%this padding in the decrypted hex dumps. The padding, we note, is the amount of
%padding neccesary, e.g. in f1 we use $0b$ bytes of padding, in $f2$ only 06.
%This is the way the PKCS5 standard reccomends it be.

\subsection{Task 5}
For this assignment we look at the error propagation of the various AES encryption
modes: ECB, CBC, CFB, and OFB. We do this by creating a text file that is at least
1000 bytes long, encrypting it and corrupting a single bit. We then decrypt
the ciphertext and explain the results. The task also asks us to
think about what will happen before we produce the results.

We encrypt the following text:\\
{\tiny \emph{
Lorem ipsum dolor sit amet, consectetur adipiscing elit. Integer lectus quam, tempor quis maximus quis, laoreet a enim. In at ligula est. Proin rhoncus aliquet varius. Nunc vulputate elit sed velit tincidunt pharetra. Nam nunc mi, vehicula ut neque nec, tincidunt maximus magna. Sed consequat iaculis risus, laoreet molestie metus ornare in. Vestibulum commodo luctus felis eget scelerisque.}\\
\emph{Etiam quis tempus massa, molestie vestibulum felis. Cras quis sollicitudin metus. Vivamus sollicitudin eros nec enim tincidunt faucibus. Maecenas laoreet ultricies convallis. Aliquam suscipit egestas neque ut efficitur. Pellentesque at neque fringilla, tincidunt urna dictum, tempus lectus. Cras a nisi dui.}\\
\emph{Aenean suscipit, massa id maximus posuere, ex risus commodo quam, sed dignissim lorem dolor at lorem. Fusce bibendum, libero nec dictum semper, risus ante maximus nisl, in tincidunt lorem metus vel mi. Proin pretium orci sit amet ornare congue. Morbi nec libero vitae diam tristique dictum dapibus nam. 
}}\\
Which rounds out to 1004 bytes.

With the commands shown below:
\begin{verbatim}
$ ciphs=(ecb cbc cfb ofb)
$ for ciph in $ciphs; do 
openssl enc -aes-128-$ciph -in document.txt -out "$ciph.bin" \ 
-k 00112233445566778899aabbccddeeff -iv 0102030405060708;
done
\end{verbatim}

To save space in this report, the files, and code to generate these 
will be included in our code.zip hand-in. 

\subsubsection{ECB}
As mentioned in the previous section, in ECB every block will be encrypted
and decrypted without regard to the previous and next blocks. Thus the
final decryption with a corrupted block will be the correct text, with only
one small part of the text corrupted.\\

\subsubsection{CBC}
In CBC, the decryption of a block, relies on the ciphertext of the previous
block. In the case of a corrupted block, the block \emph{and} its
following block, is going to be corrupted, while the rest of the blocks
are decrypted as normal.\\

\subsubsection{CFB}
CFB Suffers from the same issues, with the broken ciphertext corrupting
both the current and the next block, while letting the rest of the text
remain functional.

\subsubsection{OFB}
OFB is very comparable to cfb, with the exception that OFB carries the
encryption result through, instead of the ciphertext. Thus a broken
ciphertext only corrupts the one block, and leaves the rest of the decoded
string looking good.


\end{document}
